\section{在线方法：DPO 之后的前沿}

\subsection{什么是"在线"数据？}

\begin{figure}[H]
\centering
\includegraphics[width=0.95\textwidth]{slides-images/slide-075.jpg}
\caption{"在线"数据的定义：由策略新生成 + 由奖励模型新标注\protect\footnotemark}
\end{figure}
\footnotetext{Slides 第76页。}

"在线"在对齐语境中有两层含义：

\begin{enumerate}
    \item \textbf{数据文本的新鲜度}：数据是由当前策略模型生成的（on-policy），还是来自其他模型的历史数据（off-policy）？
    \item \textbf{标签的新鲜度}：偏好标签是训练前一次性标注的，还是在训练过程中由奖励模型动态重新标注的？
\end{enumerate}

\begin{importantbox}{PPO 与 DPO 在"在线"维度上的根本差异}
PPO \textbf{天然是在线的}：它在训练过程中不断从当前策略生成新回复，然后用奖励模型评分。DPO 则使用\textbf{固定的离线数据集}——数据中的回复来自各种不同的模型（GPT-3.5、GPT-4、Alpaca、Vicuna 等），而非当前正在训练的模型。这种数据分布的差异可能是 PPO 略优于 DPO 的根本原因。
\end{importantbox}

\subsection{在线 DPO 方法的兴起}

2024年4--5月，大量论文开始研究如何将"在线"特性引入 DPO：

\begin{figure}[H]
\centering
\includegraphics[width=0.95\textwidth]{slides-images/slide-078.jpg}
\caption{关于 online vs offline DPO 的多篇并发研究\protect\footnotemark}
\end{figure}
\footnotetext{Slides 第79页。}

主要方法包括：

\begin{itemize}
    \item \textbf{Self-Rewarding Language Models}（Meta）：让 DPO 模型自己作为 judge 来重新标注偏好数据，然后进行多轮迭代 DPO 训练
    \item \textbf{分批 DPO}：不一次性用完所有数据，而是分批训练，每批之间更新数据标签
    \item \textbf{Discriminator-Guided DPO（D2PO）}：Lambert 参与指导的项目，在 DPO 过程中动态重训奖励模型并重新标注数据
\end{itemize}

\begin{figure}[H]
\centering
\includegraphics[width=0.95\textwidth]{slides-images/slide-079.jpg}
\caption{在线 DPO 的多种实现方式\protect\footnotemark}
\end{figure}
\footnotetext{Slides 第80页。}

\subsection{D2PO：判别器引导的 DPO}

D2PO 比较了三种方案：

\begin{enumerate}
    \item[\textbf{(A)}] \textbf{标准 DPO}：固定数据集，一次性训练
    \item[\textbf{(B)}] \textbf{在线偏好优化}：使用奖励模型反复重新标注偏好数据
    \item[\textbf{(C)}] \textbf{D2PO}：不仅重新标注数据，还在训练过程中重训奖励模型
\end{enumerate}

实验使用了一个可度量的闭式任务（计算句子中名词数量作为奖励），结果表明 D2PO 比单纯重新标注偏好数据收敛更好。这说明\textbf{保持奖励模型与策略模型同步更新}是有价值的。

\subsection{工业界的迭代 RLHF}

\begin{figure}[H]
\centering
\includegraphics[width=0.95\textwidth]{slides-images/slide-080.jpg}
\caption{工业界的迭代 RLHF：Anthropic 的 Claude 和 Meta 的 Llama 2 的多轮训练过程\protect\footnotemark}
\end{figure}
\footnotetext{Slides 第81页。}

工业界的做法更加直接：

\begin{itemize}
    \item \textbf{Anthropic（Claude）}：多轮 Constitutional RL，每轮使用新的人类数据
    \item \textbf{Meta（Llama 2）}：明确说明与标注公司合作，分批收集数据，每批使用上一轮模型的检查点进行生成
\end{itemize}

\begin{knowledgebox}{工业界与学术界的"在线"差异}
工业界的"在线"意味着\textbf{不断收集新的人类数据}——每个数据点都是花钱买来的。学术界的"在线"更多是用奖励模型\textbf{重新标注已有数据}或\textbf{从策略模型生成新数据}。两者的规模和质量有本质差异：Meta 在 Llama 2 中购买了约 150 万条比较数据，远超 Chatbot Arena 收集的 80 万条。
\end{knowledgebox}

\subsection{Llama 3 的后训练策略}

Llama 3 的博客文章透露了一句意味深长的话：\textit{"Our approach to post-training is a combination of supervised fine-tuning, rejection sampling, proximal policy optimization, and direct preference optimization."}

Lambert 的解读是：Meta 在不同训练阶段使用了不同的方法——在每个检查点尝试多种方法，选择效果最好的。这是一种非常实用主义的策略：

\begin{itemize}
    \item \textbf{Rejection Sampling}（最简单）：用奖励模型对 SFT 输出排序，取最好的继续训练
    \item \textbf{DPO}（较简单）：适合快速迭代
    \item \textbf{PPO}（最复杂但最强）：在最终阶段精细调优
\end{itemize}

\begin{figure}[H]
\centering
\includegraphics[width=0.95\textwidth]{slides-images/slide-082.jpg}
\caption{Meta Llama 3 的后训练策略：结合 SFT、Rejection Sampling、PPO 和 DPO\protect\footnotemark}
\end{figure}
\footnotetext{Slides 第83页。}

\subsection{本章小结}

"在线"方法代表了对齐研究最活跃的前沿方向。核心洞察是：使用当前策略生成的数据（on-policy data）和及时更新的偏好标签，比使用固定的离线数据集能够带来更好的训练效果。在工业界，这意味着持续投入人类标注资源；在学术界，这意味着开发创新的自标注和迭代训练方法。

%% ========== Part 7: 未来方向 ==========

